在北半球局部 \(f\)-平面上，考虑均匀静水深 \(H>0\) 的不可压、无粘、自由表面浅水系统。重力加速度为 \(g>0\)，科里奥利参数 \(f>0\) 为常数；记东西、南北向深度平均速度为 \(u,v\)，自由面微小扰动为 \(\eta\)。线性化方程为
\[
u_t-fv=-g\eta_x,\qquad v_t+fu=-g\eta_y,\qquad
\eta_t+H(u_x+v_y)=0,\qquad c=\sqrt{gH}.
\]
忽略底摩擦和表面张力，能量密度均按单位水平面积、单位水密度计算。
(1) 在无限平面取沿 \(x\) 方向的傅里叶扰动 \(e^{i(kx-\omega t)}\)，其中 \(k\ne0\)。从完整的三变量线性系统求全部三个频率分支；写出两支非零频率惯性重力波的速度极化关系，并解释零频率地转平衡模。
(2) 现在令海洋占据半平面 \(y\ge0\)，海岸 \(y=0\) 为不透水刚性边界，要求 \(v(x,0,t)=0\)，且扰动随 \(y\to+\infty\) 衰减。寻找 \(v\equiv0\) 的非平凡沿岸行波，严格推导其频散关系、横向衰减长度和唯一允许的传播方向；说明改变 \(f\) 的符号会怎样改变方向。
(3) 定义相对涡度 \(\zeta=v_x-u_y\)。从原方程证明线性位涡扰动 \(q'=\zeta-f\eta/H\) 在每个固定位置守恒；再推导局域波动能量守恒律，证明科里奥利力不作功，并检查第 (2) 问的 Kelvin 波满足 \(q'=0\)。
(4) 求惯性重力波正频支的群速度，分析长波、短波及 \(f\to0^+\) 极限。对第 (2) 问取实波面
\(\eta=\eta_0e^{-y/L_R}\cos(kx-ckt)\)，计算单位海岸长度的时间平均总能量及沿岸平均能流，验证能量输运速度，并说明海岸波在 \(f\to0^+\) 时为何不再横向局域化。
